\section{Method Details}
\label{app:method-details}

This appendix specifies the execution details supporting the three stages in
Section~\ref{sec:method}. Development fixes representation and measurement rules,
independent qualification assesses them, discovery prioritizes hypotheses, and
confirmation tests the selected policies. Deduplicated task units and their
frozen near-duplicate groups do not cross these roles; instances and repeated
generations remain attached to their task unit.

\subsection{Source Representation and Scope Checks}
\label{app:representation-details}

The source records retain operation participants and outputs, requirement
targets, conditions, and composition together with exact evidence spans.
Mechanical checking permits at most one bounded replacement. One ordinary
source review then permits at most one local semantic patch before deterministic
graph compilation. The patch is restricted to affected records and their
source-supported dependencies. Compilation checks permitted types, references,
and evidence locations but cannot infer missing bindings. A positive source
review is not an independent semantic qualification.

Conditions supplying information remain distinct from the objects being
constrained and from conditions governing whether an operation executes.
For example, requiring Python does not make deletion conditional on Python;
participation in a common workflow does not spread a requirement to every
operation. Candidate construction preserves negation, thresholds, direct objects,
and compound meanings. It cannot convert an input assumption into an obligation,
select half of a compound obligation, or repair missing graph relations.

Context and scope matching are independent of the target requirement state.
The subsequent source-only assessment refers to existing operation, input, and
condition instances without adding, splitting, or rebinding facts. Missing or
ambiguous evidence remains unresolved. Unknowns can be localized only when their
source-grounded impacts are known; shared constraints or missing localization
remain blocking. Independent qualification must assess source meaning and
candidate correctness, not only graph well-formedness.

\subsection{Candidate Freeze and Intervention Execution}
\label{app:intervention-details}

Before discovery outcomes are inspected, the candidate universe, natural-state
support rules, reliability thresholds, and task-unit folds are fixed. Source
tasks must establish the relevant operation, security boundary, and observable
non-target behavior; they need not already contain experimental controls or
exact editing recipes. Missing functionality details limit the corresponding
claims without automatically invalidating an independently measurable security
comparison. Full and RD-only prioritization use identical candidate rows,
support decisions, folds, cross-fitted standardized RD scores, deterministic
tie-breaks, and fixed $K$; only Full applies the FCI adjacency criterion.

A common outcome-blind mapping resolves each method's proposal to a scoped
requirement, editing operation, and comparator. This mapping is distinct from
protocolization, which validates concrete edits and controls, and from
confirmation, which estimates effects. Each unique selected model-bound effect
is confirmed once and reused by all referring slots. Empty, failed, and
non-evaluable slots remain in the original $K$; outcomes cannot trigger
reranking, replacement, or budget reassignment.

The hypothesis freeze fixes finite editing realizations, their positive weights
summing to one, eligible population rules, arms, models, request settings,
endpoints, and inference plans. Seeded, source-stratified balanced allocation
assigns exactly one realization to each task--policy coordinate before
task-specific materialization. Only that realization's complete four-arm bundle
is produced and checked. Failure is retained without trying another realization,
replacing the task, or renormalizing weights. Realizations may consequently have
different protocolizable populations, a limitation of the mixture estimand.

Eligible pools, materialized bundles, and the complete assignment table are
frozen before generation. Within each task-instance--policy--realization--model
block, the four arms receive equal requests and randomized execution order.
Task, arm, realization, model, and request-randomness bindings are preserved.
Failed materialization is a pre-assignment record; each committed assignment
subsequently receives a terminal record and remains in its assigned arm.

\subsection{Measurement, Weighting, and Simultaneous Inference}
\label{app:inference-details}

Oracle qualification covers the admitted endpoint, languages, task families,
and code forms, using representative secure, insecure, and unknown examples
and checks for plausible differential measurement errors. Functional assessment
is arm blind and uses the complete original task. Whole-task pass requires
assessable material requirements and their fulfillment; unresolved requirements
cannot be credited as satisfied. Incomplete source specifications alone do not
determine failure. Infrastructure evidence that is missing blocks affected
inference pending resolution or replay of the same assignments under the
frozen rules.

For hypothesis $h$ and a fixed model, let $\bar Y_{ia}$ average request outcomes
for arm $a$ within each task instance and then combine instances within task
unit $i$ using frozen weights. Realization $r\in R_h$ has weight $q_{hr}>0$,
with $\sum_r q_{hr}=1$. For $n_{hr}$ eligible assigned task units allocated to
realization $r$, the arm mean is
\begin{equation}
  \widehat\mu_{ha}
  =\sum_{r\in R_h}q_{hr}
     \left(\frac{1}{n_{hr}}\sum_{i:r_{hi}=r}\bar Y_{ia}\right).
\end{equation}
Here $r_{hi}$ denotes task unit $i$'s allocated realization.
Thus realization-specific task means are combined using the original mixture
weights, not weights selected by realized failures. Within realization/source
cells, $\lambda_{hrs}=q_{hr}n_{hrs}/n_{hr}$, where $s$ indexes source strata
and $n_{hrs}$ counts task units in cell $(r,s)$. These fixed weights also apply to
arm summaries, diagnostic fractions, and sensitivity bounds. Missing required
support never creates a smaller renormalized policy.

For latent secure valid-code yield, define $L=Y$ and
$U=L+\mathbb I\{\text{valid and Oracle-unevaluable}\}$. The contrast bounds are
$\widehat\mu_{LT}-\widehat\mu_{UN}$ through
$\widehat\mu_{UT}-\widehat\mu_{LN}$ under the same weights. Oracle-unknown,
functional-unknown, and joint-unknown outcomes retain separate coverage records.
Specificity, functionality, and joint-success claims require their own
prospectively frozen secondary inference.

Let $D_{hi}=\bar Y_{iT}-\bar Y_{iN}$. The variance estimate is
$\sum_{r,s}\lambda_{hrs}^{2}s_{hrs}^{2}/n_{hrs}$, where $s_{hrs}^{2}$ is the
sample variance of $D_{hi}$ in the corresponding cell. Each bootstrap draw
samples the family-wide union of task units with replacement within each
original source stratum, retaining that stratum's original number of units.
All descendants move together, including arms, instances, overlapping effects,
and models. Each effect retains its original support, realization labels,
and $\lambda$ weights rather than refitting weights to resampled frequencies.

The primary family contains every unique selected Atomic target-minus-no-op
secure-yield model effect. A draw is invalid if any coordinate loses its frozen
minimum cell/realization support or has zero standard error. Otherwise its
statistic is the maximum across coordinates of the absolute studentized
deviation from the original estimate. Of $B$ draws, at least
$\lceil B f_{\min}\rceil$ must be valid, where $f_{\min}$ is the required
valid-draw fraction. For familywise error level $\alpha$, the critical value is the
$\lceil(1-\alpha)B_{\mathrm{valid}}\rceil$ order statistic of the valid bootstrap
maxima; intervals are the
estimate plus or minus this value times its standard error.
Failed provenance, inadequate initial support, excessive unknown coverage,
zero initial standard error, or too few valid draws makes the family
non-evaluable, with no unadjusted fallback. An empty selected union has no
eligible coordinates rather than zero-valued effects.

For a prespecified practical margin $\epsilon$, an interval wholly above
$\epsilon$ is positive-meaningful, one wholly below $-\epsilon$ is
negative-meaningful, and one contained within $[-\epsilon,\epsilon]$ is
practically null. Other evaluable intervals are inconclusive. Practical margins
and power assumptions are fixed before discovery outcomes are inspected.
The assignment-bound inference plan, including $\alpha$, support minima, draw
count, seed, and valid-draw fraction, is frozen before confirmation outcomes.

\subsection{Explanation Evidence Links}

Each explanation links its task scope and source bindings to the selected
hypothesis, editing policy, concrete arms, assignment records, measurements,
and inference output. An illustrative wording does not replace the frozen
mixture whose effect is reported. Missing evidence is explicitly unavailable;
reporting cannot alter an effect status, remove assigned tasks, or rerank
hypotheses. The explanation and action format specifies content constraints,
not an independently qualified generator. The exact adoption predicates and
their required secondary evidence remain open; this does not prevent reporting
an otherwise evaluable primary effect.
